\section{Random permutation is hard}
\label{sec:random-hard}

We prove that the cost of most permutations is $\Omega(n\log n)$ for \emph{all} BST algorithm. A similar result has been shown by Wilber~\cite{wilber} for random access sequences (that might not be permutations).
% Note that the straightforward information-theoretic argument will not work since it will only say that there exists a permutation of cost $\Omega(n \log n)$. 
More formally, we prove the following theorem. 

\begin{theorem} 
\label{thm:random-perm-hard}
There is a positive constant $c$ such that for any $L > 0$ only a fraction $1/2^L$ of the permutations $X$ have
$\opt(X) \le (\log n! - L -n -1)/c$. 
\end{theorem}   
\begin{proof}
We use Kolmogorov complexity. Observe that the shape of a tree on $k$ nodes can be encoded by $O(k)$ bits, e.g., by the preorder traversal. If the keys stored in the tree are clear from the context, the assignment of the keys to the nodes is unique. Let $T$ be the initial tree, let $X$ be any (permutation) access sequence, and let $C_X$ be the total cost of $\opt$ on $X$. We encode the shape of the initial tree with $O(n)$ bits.  Consider now an access of cost $k$. It replaces an initial segment of the current search tree, the before-tree, by another tree on the same set of nodes, the after-tree. The before-tree must include the element accessed and has size $O(k)$. We encode the before-tree and the after-tree with $O(k)$ bits and use a special marker to designate the accessed element in the before-tree. In total, we need $c \cdot (n + C_X)$ bits for the encoding for some small constant $c$. Note that $X$ can be reconstructed from the encoding. We first reconstruct the shape of the initial tree. Its labelling by keys is unique. 
Consider now any access. Since the position of the accessed element in the before-tree is marked and since the before- and the after-tree are on the same set of keys, the search tree after the access is determined. We conclude that we have an injective mapping from permutations $X$ to bitstrings of length $c \cdot (n + C_X)$. 
Assume $C_x \le (\log n! - L - n - 1)/c$. Then the encoding of $X$ has at most $\log n! - L - 1$ bits and hence $X$ is in a set of at most $n!/2^L$ permutations.
\end{proof}


